\section{Flexbox: Praktik Layout}

Setelah memahami konsep dasar Flexbox, section ini mendemonstrasikan penerapannya pada komponen UI nyata. Flexbox unggul untuk layout satu dimensi: menyusun elemen secara horizontal atau vertikal dengan distribusi ruang yang fleksibel dan alignment yang presisi. Pola-pola berikut menjadi fondasi untuk membangun hampir semua komponen antarmuka modern \cite{mdnCss}.

\begin{contoh}
\textbf{Studi Kasus: Navbar Responsif dengan Flexbox}

Navbar desktop menampilkan logo di kiri dan menu di kanan. Pada layar kecil, menu berubah menjadi vertikal.
\begin{enumerate}
  \item HTML: \code{<nav>} berisi \code{<div class="logo">} dan \code{<ul class="menu">}
  \item CSS Desktop: \code{nav \{ display: flex; justify-content: space-between; align-items: center; \}}
  \item Item \code{<ul>} juga flex: \code{.menu \{ display: flex; gap: 1rem; list-style: none; \}}
  \item Media query mobile: \code{flex-direction: column} pada \code{.menu}
\end{enumerate}
\end{contoh}

\begin{webcode}[language=CSS, caption={Navbar Flexbox responsif}]
nav {
  display: flex;
  justify-content: space-between;
  align-items: center;
  padding: 0.5rem 1rem;
  background: #2a6fb0;
}
.menu {
  display: flex;
  gap: 1rem;
  list-style: none;
  margin: 0; padding: 0;
}
.menu a { color: white; text-decoration: none; }

@media (max-width: 768px) {
  nav { flex-direction: column; }
  .menu { flex-direction: column; text-align: center; }
}
\end{webcode}

\begin{konsep}
\textbf{Pattern Flexbox yang Sering Digunakan:}
\begin{itemize}
  \item \textbf{Centering}: \code{display: flex; justify-content: center; align-items: center;} --- cara termudah untuk centering vertikal dan horizontal.
  \item \textbf{Space between}: \code{justify-content: space-between} --- logo kiri, menu kanan.
  \item \textbf{Equal columns}: setiap item \code{flex: 1} --- kolom sama rata.
  \item \textbf{Sticky footer}: body flex column + main \code{flex: 1} --- footer selalu di bawah \cite{cssTricks}.
\end{itemize}
\end{konsep}

